\cleardoublepage
\chapter{\textit{PROMPT} SISTEM}
\label{lamp:prompt-sistem}

Lampiran ini memuat seluruh \textit{prompt} yang dipakai pada implementasi sistem mitigasi halusinasi RAG-LLM. \textit{Prompt} ditampilkan dalam format kode persis seperti yang diteruskan ke model bahasa, agar reproduksibilitas eksperimen tetap terjaga. Variabel substitusi ditulis dengan kurung kurawal (misalnya \texttt{\{context\}} dan \texttt{\{question\}}) sebagaimana konvensi \textit{string formatting} Python.

\section{\textit{Prompt} Generation Jawaban Biomedis}
\label{lamp:prompt-generation}

\textit{Prompt} berikut dipakai pada tahap \textit{generation} untuk seluruh konfigurasi sistem (\textit{baseline}, +QR, +CR, +QR+CR). \textit{Prompt} dirancang agar LLM menjawab pertanyaan biomedis berbasis lima \textit{chunk} konteks yang sudah dipilih oleh \textit{retriever} dan mengakhiri jawaban dengan satu dari tiga label klasifikasi (\textit{yes}, \textit{no}, atau \textit{maybe}) pada baris terakhir. Instruksi spesifik mengenai kapan memilih label \textit{maybe} ditambahkan untuk mencegah model bersikap terlalu konservatif pada bukti yang sebenarnya cukup memadai untuk satu arah.

\begin{lstlisting}[language=,caption={\textit{Prompt} \textit{generation} jawaban biomedis (GENERATION\_PROMPT)},label={lst:generation-prompt}]
You are a medical research assistant. Answer a biomedical
yes/no/maybe question based solely on the provided scientific
abstracts.

Context from medical literature:
{context}

Question: {question}

Instructions:
- Carefully read the context and assess whether it supports
  or refutes the question.
- Provide a brief explanation (2-3 sentences) using ONLY the
  information above.
- End your response with EXACTLY ONE of these words on its
  own line: yes, no, or maybe.
  - yes   : the evidence supports the hypothesis, even if not
            perfectly conclusive
  - no    : the evidence refutes or does not support the
            hypothesis
  - maybe : ONLY if the evidence is directly contradictory
            (some findings say yes, others say no), or if the
            context contains no relevant information at all
- IMPORTANT: If the evidence leans in one direction, even
  partially, choose yes or no. Do NOT use maybe simply
  because the evidence is limited or not 100% certain.

Answer:
\end{lstlisting}

Variabel \texttt{\{context\}} pada \textit{Listing}~\ref{lst:generation-prompt} diisi dengan lima \textit{chunk} hasil \textit{retrieval} yang sudah disusun secara berurutan dan diberi penanda nomor serta label \textit{section} sumber paper. Format penyusunan konteks dilakukan oleh fungsi pembentuk konteks pada modul \textit{generation}, sebagaimana ditunjukkan pada \textit{Listing}~\ref{lst:context-format}.

\begin{lstlisting}[language=Python,caption={Fungsi pembentuk \texttt{\{context\}} untuk diteruskan ke \textit{prompt}},label={lst:context-format}]
def generate_answer(query, retrieved):
    context = "\n\n".join(
        f"[{i}] ({r.document.section_label}): {r.document.text}"
        for i, r in enumerate(retrieved, 1)
    )
    return openai_generate(
        GENERATION_PROMPT.format(context=context, question=query),
        max_tokens=300, temperature=TEMPERATURE
    )
\end{lstlisting}

\section{\textit{Prompt} Query Rewriting}
\label{lamp:prompt-qr}

\textit{Prompt} berikut dipakai pada modul \textit{Query Rewriting} (QR) yang berjalan sebelum tahap \textit{retrieval} pada konfigurasi +QR dan +QR+CR. Tugasnya sederhana, yaitu menulis ulang pertanyaan pengguna menjadi satu pertanyaan baru yang lebih spesifik dan lebih kaya istilah biomedis. \textit{Prompt} ini berisi empat aturan utama, yaitu (1) membuat pertanyaan lebih spesifik dengan menambahkan istilah medis yang relevan, (2) mengekspansi akronim seperti \texttt{MI} menjadi \textit{myocardial infarction}, (3) mempertahankan jenis pertanyaan agar tetap dapat dijawab dengan \textit{yes}, \textit{no}, atau \textit{maybe}, serta (4) mengeluarkan hanya pertanyaan baru tanpa penjelasan tambahan.

\begin{lstlisting}[language=,caption={\textit{Prompt} \textit{Query Rewriting} dengan pendekatan \textit{single-query rewriting}},label={lst:qr-prompt}]
You are a query rewriting assistant for a biomedical
question-answering system.
Rewrite the following medical question to improve retrieval
from a PubMed research database.

Rules:
1. Be more specific and add relevant medical/scientific
   terminology.
2. Expand abbreviations (e.g. "MI" -> "myocardial infarction").
3. Preserve the original yes/no/maybe answerable intent.
4. Output ONLY the rewritten question, no explanations.

Original question: {query}

Rewritten question:
\end{lstlisting}

Keluaran LLM langsung dipakai sebagai kueri baru pada tahap \textit{retrieval}, baik BM25 maupun \textit{dense}. Apabila pertanyaan hasil tulis-ulang terlalu pendek (kurang dari 10 karakter) atau pemanggilan LLM gagal, sistem akan memakai pertanyaan asli sebagai \textit{fallback} agar \textit{pipeline} tetap berjalan tanpa harus menghentikan proses. \textit{Listing}~\ref{lst:qr-function} memperlihatkan implementasi fungsi yang menjalankan pendekatan ini.

\begin{lstlisting}[language=Python,caption={Fungsi \textit{rewriter} \textit{single-query} (\texttt{rewrite\_query})},label={lst:qr-function}]
def rewrite_query(query: str) -> str:
    """Reformulasi query menjadi satu pertanyaan baru yang
    lebih spesifik dan kaya istilah biomedis."""
    prompt = QUERY_REWRITE_PROMPT.format(query=query)
    try:
        rewritten = openai_generate(
            prompt, max_tokens=150, temperature=0.3
        )
        rewritten = rewritten.strip().replace('\n', ' ')
        return rewritten if len(rewritten) >= 10 else query
    except Exception:
        return query  # fallback ke pertanyaan asli
\end{lstlisting}

\section{Parameter Pemanggilan Model Bahasa}
\label{lamp:parameter-llm}

Seluruh pemanggilan model bahasa pada tahap \textit{generation} dan \textit{Query Rewriting} memakai parameter yang ditetapkan sama untuk seluruh konfigurasi agar perbandingan antarmodel bersifat \textit{apple-to-apple}. Ringkasan parameter tersebut disajikan pada Tabel~\ref{tbl:parameter-llm}.

\begin{table}[H]
\centering
\caption{Parameter Pemanggilan Model Bahasa pada Tahap \textit{Generation} dan \textit{Query Rewriting}}
\label{tbl:parameter-llm}
\begin{tabular}{|l|c|c|}
\hline
\textbf{Parameter} & \textbf{Generation} & \textbf{Query Rewriting} \\ \hline
\textit{Temperature}            & 0{,}0 & 0{,}3 \\ \hline
\textit{Max tokens}             & 300   & 150   \\ \hline
\textit{Seed}                   & 42    & 42    \\ \hline
\end{tabular}
\end{table}

Penggunaan \textit{temperature} 0{,}0 pada tahap \textit{generation} menjamin determinisme jawaban sehingga eksperimen yang sama menghasilkan keluaran identik antar-pemanggilan. Untuk \textit{Query Rewriting}, \textit{temperature} dinaikkan menjadi 0{,}3 agar LLM dapat menghasilkan reformulasi yang lebih kaya istilah biomedis sambil tetap menjaga makna pertanyaan asli.
